% 成对项的分解（引理 lem:graph-setfn-pairwise-decomposition），两变量例子：
% phi_{1,2}(0,0)=0, (0,1)=3, (1,0)=5, (1,1)=0；w=(3+5-0-0)/2=4，alpha_1=5-0-4=1，alpha_2=3-0-4=-1。
% 逐格核对：(0,0) 0+0+0+0=0；(0,1) 0+0-1+4=3；(1,0) 0+1+0+4=5；(1,1) 0+1-1+0=0。
\begin{tikzpicture}[x=1mm,y=1mm,line width=.6pt,font=\figurefont\fontsize{8}{10}\selectfont]
  \def\cs{5.6}
  % #1 x 偏移，#2-#5 四格数值 (00,01,10,11)，#6 填色(四格)，#7 标题
  \newcommand{\ptable}[8]{%
    \begin{scope}[xshift=#1 mm]
      \foreach \r/\c/\v/\f in {0/0/#2/#6,0/1/#3/#7,1/0/#4/#7,1/1/#5/#6}{
        \filldraw[fill=\f,draw=FigureStroke] (\c*\cs,-\r*\cs) rectangle (\c*\cs+\cs,-\r*\cs-\cs);
        \node[FigureInk] at (\c*\cs+\cs/2,-\r*\cs-\cs/2) {$\v$};
      }
      \foreach \c in {0,1}{\node[FigureMuted,font=\figurefont\fontsize{7}{8}\selectfont] at (\c*\cs+\cs/2,1.9) {$\c$};}
      \foreach \r in {0,1}{\node[FigureMuted,font=\figurefont\fontsize{7}{8}\selectfont] at (-1.6,-\r*\cs-\cs/2) {$\r$};}
      \node[FigureInk,align=center,anchor=north,font=\figurefont\fontsize{7.5}{10.5}\selectfont] at (\cs,-2*\cs-1.4) {#8};
    \end{scope}}
  \ptable{0}{0}{3}{5}{0}{FigureBlueFill}{FigureBlueFill}{$\phi_{1,2}(x_1,x_2)$}
  \node[FigureInk] at (14.7,-\cs) {$=$};
  \ptable{20}{0}{0}{0}{0}{white}{white}{常数\\$p_{0,0}=0$}
  \node[FigureInk] at (34.7,-\cs) {$+$};
  \ptable{40}{0}{0}{1}{1}{white}{white}{$\alpha_1x_1$\\$\alpha_1=1$}
  \node[FigureInk] at (54.7,-\cs) {$+$};
  \ptable{60}{0}{-1}{0}{-1}{white}{white}{$\alpha_2x_2$\\$\alpha_2=-1$}
  \node[FigureInk] at (74.7,-\cs) {$+$};
  \ptable{80}{0}{4}{4}{0}{white}{FigureRedFill}{$w_{1,2}\,|x_1-x_2|$\\$w_{1,2}=4$}
  % 行列变量说明
  \node[FigureMuted,anchor=east,font=\figurefont\fontsize{7}{8}\selectfont] at (-2.6,-\cs) {$x_1$};
  \node[FigureMuted,anchor=south,font=\figurefont\fontsize{7}{8}\selectfont] at (\cs,3.2) {$x_2$};
\end{tikzpicture}
